\section{Truncated spectral-row energy}
\label{restore:tail-section}

The next lemma controls large spectral rows in expectation. Its proof
uses a recurrence at a single truncation threshold, together with
concentration of independent deleted-row innovations.

Work under the fixed structural assumptions of
\cref{prop:sparse-floor}, with $np\to\infty$ and $np^2\le C_1$.
For the high-degree cap in \cref{restore:algorithm}, recall
\[
 \tau=b_0\bar D,\qquad
 G_j=\max(D_j,\tau),\qquad
 W=\tau^{1/2}G^{-1/2},\qquad \widetilde A=WAW.
\]
Let $U$ be an orthonormal basis of the informative $K$-dimensional
eigenspace and write $\widetilde AU=U\Lambda$.
All bounds below are uniform in the fixed compact structural class.
In particular,
\begin{equation}
 0\le{P_{z_i z_j}}\le Cp,\qquad
 \norm{{(ZPZ^\top)_{i\cdot}}}_2\le C{np}/\sqrt n.
 \label{restore:tail-population-rows}
\end{equation}

Fix a probability parameter $r>0$.  Apply
\cref{prop:sparse-floor} with a sufficiently larger fixed parameter to
the full graph and all graphs obtained by deleting at most two vertices.
Let $\mathcal E_r$ be the intersection of the resulting geometry and
two-sided total-degree events.  Restricting the population matrix
preserves the fixed structural bounds and comparable block masses.
Thus
\begin{equation}
 \PP(\mathcal E_r^c)\le C_rn^{-r}.
 \label{restore:tail-geometry-probability}
\end{equation}
On this event, for every graph in the deletion family, its recomputed
threshold is comparable to ${np}$, its total degree is at most $C_r{n^2p}$,
its centered weighted noise has norm at most $C_r\sqrt {np}$, its
informative eigenvalues have absolute values at least $c_r{np}$, and its
remaining eigenvalues have absolute values at most $C_r\sqrt {np}$.
For a deleted training graph, use the analogous event formed from
that graph and its further deletions. This event is measurable from
the training graph alone and has the same probability bound.

The following deterministic properties of the cap will be useful:
\begin{equation}
 W_j\le1,\qquad W_j^2D_j\le\tau,
 \qquad \norm{\widetilde A}_{\op}\le\tau.
 \label{restore:tail-floor-operator}
\end{equation}
Indeed, $G\ge\diag(A\one)$ and the inequality
$2|ab|\le a^2+b^2$ gives
$|x^\top Qx|\le\sum_j(D_j/G_j)x_j^2\le\norm x_2^2$,
where $Q=\widetilde A/\tau$.

\begin{lemma}[Truncated energy for the capped-adjacency eigenvectors]
\label{restore:tail-energy-lemma}
There are fixed constants $R_0,C,c>0$ such that, for every fixed
$R\ge R_0$ and all sufficiently large $n$,
\begin{equation}
 \E\!\left[\ind{\mathcal E_r}
       \sum_{j=1}^n\norm{U_{j\cdot}}_2^2
                  \ind{\sqrt n\norm{U_{j\cdot}}_2>R}\right]
 \le CR^2e^{-c{np}\log R}.
 \label{restore:tail-energy}
\end{equation}
The constants $R_0,C,c$ depend on the fixed structural bounds,
$b_0$, and $r$, but not on $R$ or $n$.  The sufficiently-large-$n$
condition may depend on the chosen fixed $R$.  The same result holds
for each recomputed deleted graph with its training-measurable
geometry event.
\end{lemma}

\begin{proof}
We give the weight perturbation, the same-threshold energy recurrence,
and its expectation bounds separately.

\medskip\noindent\textbf{Recomputed weights and the eigenspace residual.}
For $J=[n]\setminus\{i\}$, let $W^{(-i)}$ and $U^{(-i)}$ be
computed from $A_{JJ}$, with $U^{(-i)}$ embedded by a zero row at $i$.
Outside the geometry event choose an arbitrary deterministic
orthonormal fallback depending only on $A_{JJ}$.
For an all-zero training graph set $W^{(-i)}=I$ and use that
orthonormal fallback. Thus the proof-only completion always has
$\norm{W^{(-i)}}_{\op}\le1$, and the independent innovation below
is defined for every training graph.
Put $(\Gamma_i)_{jj}=W_j^{(-i)}/W_j-1$ for $j\in J$ and
$(\Gamma_i)_{ii}=0$.
The exact threshold difference is
\[
 \tau^{(-i)}-\tau
 =b_0\frac{\bar D-2D_i}{n-1}.
\]
Comparing the two square-root caps, using the comparability of the
thresholds to ${np}$, gives on $\mathcal E_r$
\begin{equation}
 |(\Gamma_i)_{jj}|
 \le C\left\{\frac{A_{ij}}{G_j}
             +\frac{{np}+D_i}{{n^2p}}\right\},
 \qquad \norm{\Gamma_i}_{\op}\le C/(np).
 \label{restore:tail-weight-change}
\end{equation}
This also covers a crossing of the cap threshold. It follows either
by splitting the cases in the two maxima or by comparing the logarithms
of $\sqrt{\tau/\max(D_j,\tau)}$; the relative threshold change is
$O(({np}+D_i)/({n^2p}))$.  The last bound uses $D_i\le n$ and ${np}\le Cn$.

There is an orthogonal matrix $O_i$ such that
\begin{equation}
 \norm{U-U^{(-i)}O_i}_F
 \le C\norm{U_{i\cdot}}_2+C\norm{\Gamma_iU}_F.
 \label{restore:tail-loo-residual}
\end{equation}
To verify this, write
$\widetilde A=\left(\begin{smallmatrix}0&m_i^\top\\m_i&B_i\end{smallmatrix}\right)$.
By \eqref{restore:tail-floor-operator},
$\norm{m_i}_2\le\sqrt\tau$ and $\norm{B_i}_{\op}\le\tau$.
Also $B_iU_J=U_J\Lambda-m_iU_{i\cdot}$.
The lower block of the recomputed deleted matrix is
$(I+\Gamma_i)B_i(I+\Gamma_i)$.
Subtracting the full matrix and applying the difference to $U$ gives
a Frobenius residual bounded by
$C{np}(\norm{U_{i\cdot}}_2+\norm{\Gamma_iU}_F)$.
The full informative spectrum is separated from the deleted
noninformative spectrum by $c{np}$.
The complementary-space Sylvester equation and orthogonal Procrustes
alignment give \eqref{restore:tail-loo-residual}.
Only separation from the complementary spectrum is used; the
informative eigenvalues can repeat or have either sign.

\medskip\noindent\textbf{Truncation and summable weight changes.}
Set $t=R/\sqrt n$.  Let $f_t$ be radial projection onto the Euclidean
ball of radius $t$, applied to each row, and put
\[
 \bar U=f_t(U),\qquad E_U=U-\bar U,\qquad
 \mathcal T_R(U)=\sum_j\norm{U_{j\cdot}}_2^2
                         \ind{\norm{U_{j\cdot}}_2>t}.
\]
The radial projection is $1$-Lipschitz and commutes with orthogonal
rotations.  Moreover $\norm{E_U}_F^2\le \mathcal T_R(U)$,
$\max_j\norm{\bar U_{j\cdot}}_2\le t$, and $\norm{\bar U}_F\le\sqrt K$.
Define
\[
 \xi_i=W_i\left(\sqrt{D_i}+C{np}/\sqrt n\right).
\]
By \eqref{restore:tail-floor-operator}, $\xi_i^2\le C{np}$.
Equation \eqref{restore:tail-weight-change} yields
\[
 \norm{\Gamma_i\bar U}_F
 \le Ct\frac{\sqrt{D_i}}{np}+C\frac{{np}+D_i}{{n^2p}},
\]
and therefore
\begin{equation}
 \frac{C\xi_i}{np}\norm{\Gamma_i\bar U}_F
 \le Ct\left\{\frac1{np}+\frac{\sqrt{D_i}}{{(np)}^{3/2}}\right\}.
 \label{restore:tail-bounded-weight}
\end{equation}
For the global term, use $t\ge n^{-1/2}$, $D_i\le n$, and
${np}\le Cn$: $D_i/(n{(np)}^{3/2})\le t\sqrt{D_i}/{(np)}^{3/2}$ and
$1/(n\sqrt {np})\le Ct/(np)$.
The square-root cap does not give $W_iD_i\le C{np}$. The
right side of \eqref{restore:tail-bounded-weight} can exceed a
sufficiently small fixed multiple of $t$, for large ${np}$, only if
$D_i\ge c_*{(np)}^3$.

The unbounded part satisfies the summable estimate
\begin{equation}
 \sum_i\xi_i^2\norm{\Gamma_iE_U}_F^2\le C\norm{E_U}_F^2.
 \label{restore:tail-summable-weight}
\end{equation}
Indeed, after interchanging $i,j$, the local coefficient of
$\norm{(E_U)_{j\cdot}}_2^2$ is at most
$C{np}D_j/G_j^2\le C$.
The global coefficient is bounded by
\[
 \frac{C}{n^2{np}}\sum_i({np}+D_i)^2\le C,
\]
because $\sum_iD_i\le C{n^2p}$ and
$\sum_iD_i^2\le n\sum_iD_i$.
These estimates do not use a maximum-degree event.

\medskip\noindent\textbf{Independent innovations and the recurrence.}
Define the unrotated, training-measurable-coefficient innovation
\begin{equation}
 K_i=(A_{iJ}-{(ZPZ^\top)_{iJ}})W^{(-i)}f_t(U_J^{(-i)}).
 \label{restore:tail-innovation}
\end{equation}
Let $E_0=W(A-{ZPZ^\top})W$.
The eigenvector equation is
\[
 U=W{ZPZ^\top} WU\Lambda^{-1}+E_0\bar U\Lambda^{-1}+E_0E_U\Lambda^{-1}.
\]
The first term has every row bounded by $C/\sqrt n$, by
\eqref{restore:tail-population-rows}.
For the second, compare $W_J\bar U_J$ to
$W^{(-i)}f_t(U_J^{(-i)})O_i$ and use radial Lipschitzness,
\eqref{restore:tail-loo-residual}, and
\eqref{restore:tail-weight-change}.  This gives
\begin{align}
 \norm{U_{i\cdot}}_2\le{}&
 C/\sqrt n+(C/(np))\norm{K_i}_2
 +(C\xi_i/(np))\norm{U_{i\cdot}}_2\nonumber\\
 &+(C\xi_i/(np))\norm{\Gamma_i\bar U}_F
 +(C\xi_i/(np))\norm{\Gamma_iE_U}_F
 +(C/(np))\norm{(E_0E_U)_{i\cdot}}_2.
 \label{restore:tail-row-recursion}
\end{align}
The completed-mean diagonal contributes at most $Cpt/(np)\le Ct/n$
and is absorbed in the bounded terms; it is not treated as an
observed self-loop.
Although $O_i$ may depend on the tested row, the argument only uses
$\norm{K_iO_i}_2=\norm{K_i}_2$.
No conditional concentration is applied to rotated coordinates.

Take $R$ sufficiently large and then ${np}$ sufficiently large to absorb
the baseline and self terms.  Retain from
\eqref{restore:tail-bounded-weight} the excess
$Ct\sqrt{D_i/{(np)}^3}\ind{D_i\ge c_*{(np)}^3}$.
For $\norm{U_{i\cdot}}_2>t$, separate the cases
$\norm{K_i}_2>c'{np}t$ and $\norm{K_i}_2\le c'{np}t$ and square.
The resulting bound is
\begin{align}
 \norm{U_{i\cdot}}_2^2\ind{\norm{U_{i\cdot}}_2>t}
 \le{}&\frac C{{(np)}^2}\norm{K_i}_2^2
                         \ind{\norm{K_i}_2>c'{np}t}\nonumber\\
 &+\frac C{{(np)}^2}\norm{(E_0E_U)_{i\cdot}}_2^2
 +\frac{C\xi_i^2}{{(np)}^2}\norm{\Gamma_iE_U}_F^2
 +\frac{Ct^2}{{(np)}^3}D_i\ind{D_i\ge c_*{(np)}^3}.
 \label{restore:tail-pointwise-recursion}
\end{align}
Summing, using $\norm{E_0}_{\op}^2\le C{np}$ and
\eqref{restore:tail-summable-weight}, and multiplying by
$\ind{\mathcal E_r}$ gives
\begin{align}
 \mathcal T_R(U)\ind{\mathcal E_r}
 \le{}&\frac C{{(np)}^2}\sum_i\norm{K_i}_2^2
                         \ind{\norm{K_i}_2>c'{np}t}
 +\frac C{np}\mathcal T_R(U)\ind{\mathcal E_r}\nonumber\\
 &+\frac{Ct^2}{{(np)}^3}\sum_iD_i\ind{D_i\ge c_*{(np)}^3}.
 \label{restore:tail-energy-recursion}
\end{align}
Both energy terms use the same threshold $R$.

\medskip\noindent\textbf{Expectation bounds.}
Conditional on $A_{JJ}$, \eqref{restore:tail-innovation} is a sum
of independent centered vectors.  In each coordinate its summands
are bounded by $t$, and its total variance is at most
\[
 Cp\norm{W^{(-i)}f_t(U^{(-i)})}_F^2\le Cp.
\]
These bounds hold for every training graph, without conditioning on
$\mathcal E_r$.  Bennett's inequality in the fixed $K$ coordinates,
followed by integration of the squared-norm tail, gives
\begin{equation}
 \E\!\left[\norm{K_i}_2^2
                        \ind{\norm{K_i}_2>c'{np}t}\right]
 \le C{(np)}^2t^2e^{-c_1{np}\log R}.
 \label{restore:tail-innovation-moment}
\end{equation}
To see the exponent, set $v=Cp$ and
$h(u)=(1+u)\log(1+u)-u$.
At $x=c'{np}t$, the Bennett exponent
$(v/t^2)h(tx/v)$ is at least $c_1{np}\log R$, since
$tx/v$ is proportional to $R^2$.
Beyond that point the derivative of the exponent is at least
$c\log R/t$, which gives the same bound upon tail integration.

Every $D_i$ is a Bernoulli sum with mean at most $C{np}$.
Its Chernoff bound and tail integration yield
\begin{equation}
 \E[D_i\ind{D_i\ge c_*{(np)}^3}]
 \le C{(np)}^3e^{-c_2{(np)}^3\log {np}}.
 \label{restore:tail-large-degree-moment}
\end{equation}
The left side is zero if $c_*{(np)}^3>n$.
Taking expectations in \eqref{restore:tail-energy-recursion},
absorbing $C/(np)$, and using $nt^2=R^2$ now gives
\[
 \E[\mathcal T_R(U)\ind{\mathcal E_r}]
 \le CR^2\{e^{-c_1{np}\log R}+e^{-c_2{(np)}^3\log {np}}\}.
\]
For each fixed $R$, the second exponential is eventually bounded by
the first after reducing $c_1$.  This proves
\eqref{restore:tail-energy}.  The proof is uniform under deletion
of a fixed number of vertices and uses their recomputed weights,
so it also proves the asserted training-graph version.
The same proof applies to any event supplying the stated aggregate
geometry, including the exponentially accurate event of
\cref{restore:high-aggregate}. The innovation and degree-tail bounds
are unconditional before restriction to that event.
\end{proof}

\begin{corollary}[Exceptional coefficients are hit with exponentially small mass]
\label{restore:tail-hit-corollary}
On a training graph with the preceding cap, put
$H=G^{-1/2}UU^\top G^{1/2}$.
Let $f$ be any training-measurable vector satisfying
$\norm f_\infty\le C_f$, and let $\varepsilon=(H-I)f$.
For every prescribed fixed $s_0>0$, there is a fixed $L_\varepsilon<\infty$
such that, for any independent tested Bernoulli row with probabilities
$p_j\le Cp$,
\begin{equation}
 \E\!\left[\ind{\mathcal E_r}
                 \sum_{j:|\varepsilon_j|>L_\varepsilon}p_j\right]
 \le C_{s_0}e^{-s_0{np}}.
 \label{restore:tail-hit-mass}
\end{equation}
Here $\mathcal E_r$ is the training-measurable geometry event.
\end{corollary}

\begin{proof}
On that event, $G_j\ge c{np}$ and $\sum_jG_j\le C{n^2p}$.
For a fixed constant $C$ determined by these two bounds,
\[
 |\varepsilon_j|
 \le C_f+\frac{\norm{U_{j\cdot}}_2}{\sqrt{G_j}}
                  \norm{U^\top G^{1/2}f}_2
 \le C_f+CC_f\sqrt n\norm{U_{j\cdot}}_2.
\]
For $L_\varepsilon=C_f+CC_fR$, the event
$|\varepsilon_j|>L_\varepsilon$ implies $\sqrt n\norm{U_{j\cdot}}_2>R$.
Thus \cref{restore:tail-energy-lemma} gives
\begin{equation}
 \E\!\left[\ind{\mathcal E_r}
                    \sum_j\ind{|\varepsilon_j|>L_\varepsilon}\right]
 \le\frac n{R^2}\E[\mathcal T_R(U)\ind{\mathcal E_r}]
 \le Cne^{-c{np}\log R}.
 \label{restore:tail-coefficient-count}
\end{equation}
Multiplication by $Cp$ bounds the expected hit mass by
$C{np} e^{-c{np}\log R}$.
Choose a fixed $R$ with $c\log R>s_0+1$ and absorb the factor ${np}$.
Independence is required only to interpret the conditional probability
that the tested row hits the training exceptional set; no independence
between the training geometry event and its eigenvectors is asserted.
\end{proof}

\begin{corollary}[Average exceptional target-row mass]
\label{restore:tail-target-corollary}
For every fixed degree cutoff $L_D$ and every fixed $s_0>0$, there is
a fixed $L_0$ such that
\begin{equation}
 \frac1n\sum_i\PP\!\left\{
    \sqrt{G_i}\norm{U_{i\cdot}}_2>L_0\sqrt{p},\quad
    D_i\le L_D{np},\quad \mathcal E_r\right\}
 \le C_{s_0,L_D}e^{-s_0{np}}.
 \label{restore:tail-target-mass}
\end{equation}
\end{corollary}

\begin{proof}
On the displayed event, $G_i\le C_{L_D}{np}$.
Choose $L_0=R\sqrt{C_{L_D}}$.
The target-row exception then implies
$\sqrt n\norm{U_{i\cdot}}_2>R$.
Apply \eqref{restore:tail-energy} and divide its corresponding
expected count by $n$.  Choosing fixed $R$ sufficiently large gives
\eqref{restore:tail-target-mass}.
\end{proof}

